% Part: sets-functions-relations
% Chapter: functions
% Section: kinds

\documentclass[../../../include/open-logic-section]{subfiles}

\begin{document}

\olfileid[fa-IR]{sfr}{fun}{kin}
\olsection{انواع توابع}

\begin{explain}
مفید است که نوعی رده‌بندی برای برخی از
انواع توابعی که بیشتر با آن‌ها روبه‌رو می‌شویم معرفی کنیم.

برای آغاز، ممکن است بخواهیم توابعی را بررسی کنیم که این ویژگی را دارند
که هر عضو هم‌دامنه مقداری از تابع است. چنین
توابعی !!{surjective} نامیده می‌شوند و می‌توان آن‌ها را چنان‌که در
\olref{fig:surjective} آمده است به تصویر کشید.

\begin{figure}
  \olasset{assets/diagrams/surjective.tikz}
  \caption{تابع !!^a{surjective} هر !!{element} از
    هم‌دامنه را به‌عنوان یک مقدار دارد.}
  \ollabel{fig:surjective}
\end{figure}
\end{explain}

\begin{defn}[تابع !!^{surjective}]
تابع $f \colon A \rightarrow B$ اگر و تنها اگر $B$
بردِ~$f$ نیز باشد، \emph{!!{surjective}} است؛ یعنی برای هر $y \in B$ دست‌کم
یک $x \in A$ وجود دارد چنان‌که~$f(x) = y$، یا به زبان نمادین:
\[
  (\forall y \in B)(\exists x \in A)f(x) = y.
\]
چنین تابعی را !!a{surjection} از $A$ به $B$ می‌نامیم.
\end{defn}

\begin{explain}
اگر می‌خواهید نشان دهید که $f$ !!a{surjection} است، باید نشان دهید
که هر شیء در هم‌دامنهٔ $f$ مقدار $f(x)$ برای یک
ورودی $x$ است.

توجه کنید که هر تابعی !!a{surjection} \emph{القا می‌کند}. چراکه
با داشتن تابع $f \colon A \to B$، فرض کنید $f' \colon A \to \ran{f}$ به‌وسیلهٔ
$f'(x) = f(x)$ تعریف شود. چون $\ran{f}$ \emph{بنا بر تعریف} برابر است با
$\Setabs{f(x) \in B}{x \in A}$، تضمین می‌شود که این تابع $f'$
!!a{surjection} باشد
\end{explain}

\begin{explain}
حال هر تابع هر ورودی ممکن را به خروجی یکتایی نگاشت می‌کند. اما
توابعی نیز وجود دارند که هرگز ورودی‌های متفاوت را به
خروجی یکسان نگاشت نمی‌کنند. چنین توابعی !!{injective} نامیده می‌شوند و می‌توان آن‌ها را
چنان‌که در \olref{fig:injective} آمده است به تصویر کشید.
\begin{figure}
  \olasset{assets/diagrams/injective.tikz}
  \caption{تابع !!^a{injective} هرگز دو
    آرگومان متفاوت را به مقدار یکسان نگاشت نمی‌کند.}
  \ollabel{fig:injective}
\end{figure}
\end{explain}

\begin{defn}[تابع !!^{injective}]
تابع $f \colon A \rightarrow B$ \emph{!!{injective}} است اگر و تنها اگر برای
هر $y \in B$ حداکثر یک $x \in A$ وجود داشته باشد چنان‌که~$f(x) = y$.
چنین تابعی را !!a{injection} از $A$ به~$B$ می‌نامیم.
\end{defn}

\begin{explain}
اگر می‌خواهید نشان دهید که $f$ !!a{injection} است، باید نشان دهید که
برای هر $x$ و $y$ که !!{element}sی از دامنهٔ $f$ هستند، اگر $f(x)=f(y)$، آنگاه
$x=y$.
\end{explain}

\begin{ex}
تابع ثابت $f\colon \Nat \to \Nat$ که با $f(x) = 1$ داده می‌شود،
نه !!{injective} است و نه !!{surjective}.

تابع همانی $f\colon \Nat \to \Nat$ که با $f(x) = x$ داده می‌شود،
هم !!{injective} است و هم !!{surjective}.

تابع پسین $f \colon \Nat \to \Nat$ که با $f(x) = x+1$ داده می‌شود،
!!{injective} است اما !!{surjective} نیست.

تابع $f \colon \Nat \to \Nat$ که به‌صورت زیر تعریف می‌شود:
\[
  f(x) =
  \begin{cases}
    \frac{x}{2} & \text{اگر $x$ زوج باشد} \\
    \frac{x+1}{2} & \text{اگر $x$ فرد باشد.}
  \end{cases}
\]
!!{surjective} است اما !!{injective} نیست.
\end{ex}

\begin{explain}
اغلب می‌خواهیم توابعی را بررسی کنیم که هم
!!{injective} و هم !!{surjective} هستند. چنین توابعی را
!!{bijective} می‌نامیم. آن‌ها مانند تابعی هستند که در
\olref{fig:bijective} به تصویر کشیده شده است. !!^{bijection}s را گاهی
\emph{تناظرهای یک‌به‌یک} نیز می‌نامند، زیرا اعضای
هم‌دامنه را به‌طور یکتا با اعضای دامنه جفت می‌کنند.
\begin{figure}
  \olasset{assets/diagrams/bijective.tikz}
  \caption{تابع !!^a{bijective} اعضای
    هم‌دامنه را به‌طور یکتا با اعضای دامنه جفت می‌کند.}
  \ollabel{fig:bijective}
\end{figure}
\end{explain}

\begin{defn}[!!^{bijection}]
تابع $f \colon A \to B$ \emph{!!{bijective}} است اگر و تنها اگر هم
!!{surjective} و هم !!{injective} باشد. چنین تابعی را
!!a{bijection} از $A$ به~$B$ (یا میان $A$ و~$B$) می‌نامیم.
\end{defn}

\end{document}
